\subsection{MONOID OF FRACTIONS OF A COMMUTATIVE MONOID}

In this no., e will denote the identity element of a monoid and \(x^{*}\) the inverse of an invertible element x of E.

Let E be a commutative monoid, S a subset of E and S' the submonoid of E generated by S.

Lemma 1. In \(\mathbf{E} \times \mathbf{S}'\) the relation \(\mathbf{R}_{\xi}^{x}, y_{\xi}^{x}\) defined by:

``there exist \(a, b \in E\) and \(p, q, s \in S'\) such that \(x = (a, p)\) , \(y = (b, q)\) and \(aqs = bps\) ''

is an equivalence relation compatible with the law on the product monoid \(\mathbf{E} \times \mathbf{S}'\) .

It is immediate that R is reflexive and symmetric. Let \(x = (a, p)\) , \(y = (b, q)\) and \(z = (c, r)\) be elements of \(E \times S'\) such that \(R\xi x, y\xi\) and \(R\xi y, z\xi\) hold. Then there exist two elements s and t of \(S'\) such that

\[
a q s = b p s, \quad b r t = c q t,
\]

whence it follows that

\[
a r (s t q) = b p s r t = c p (s t q)
\]

and hence \(\mathbf{R}\{x, z\}\) holds, for \(stq\) belongs to \(\mathbf{S}'\) . The relation \(\mathbf{R}\) is therefore transitive.

Further, let \(x = (a, p), y = (b, q), x' = (a', p')\) , and \(y' = (b', q')\) be elements of \(E \times S'\) such that \(R\xi x, y\xi\) and \(R\xi x', y\xi\) hold. There exist \(s\) and \(s'\) in \(S'\) such that

\[
a q s = b p s, \quad a ^ {\prime} q ^ {\prime} s ^ {\prime} = b ^ {\prime} p ^ {\prime} s ^ {\prime},
\]

whence it follows that \((aa')(qq')(ss') = (bb')(pp')(ss')\) and hence \(\mathbf{R}\{\chi x', yy'\}\) for \(ss' \in S'\) . The equivalence relation \(\mathbf{R}\) is therefore compatible with the law of composition on \(\mathbf{E} \times \mathbf{S}'\) .

The quotient magma \((\mathbf{E} \times \mathbf{S}')/\mathbf{R}\) is a commutative monoid.

DEFINITION 7. Let E be a commutative monoid, S a subset of E and S' the submonoid of E generated by S. The quotient monoid (E × S')/R, where the equivalence relation R is as described in Lemma 1, is denoted by E \(_{S}\) and is called the monoid of fractions† of E associated with S (or with denominators in S).

For \(a \in \mathbf{E}\) and \(p \in \mathbf{S}'\) the class of \((a, p)\) modulo \(\mathbf{R}\) is in general denoted by \(a/p\) and called the fraction with numerator \(a\) and denominator \(p\) . Then by definition \((a/p)\) . \((a'/p') = aa'/pp'\) . The fractions \(a/p\) and \(a'/p'\) are equal if and only if there exists \(s\) in \(\mathbf{S}'\) with \(spa' = sp'a\) ; if so, there exist \(\sigma\) and \(\sigma'\) in \(\mathbf{S}'\) with \(a\sigma = a'\sigma'\) and \(p\sigma = p'\sigma'\) . In particular, \(a/p = sa/sp\) for \(a \in \mathbf{E}\) and \(s, p\) in \(\mathbf{S}'\) . The identity element of \(E_S\) is the fraction \(e/e\) .

Let \(a / e = \varepsilon(a)\) for all \(a \in E\) . The above shows that \(\varepsilon\) is a homomorphism of \(E\) into \(E_S\) , called canonical. For all \(p \in S'\) , \((p / e) \cdot (e / p) = e / e\) and hence \(e / p\) is the inverse of \(\varepsilon(p) = p / e\) ; every element of \(\varepsilon(S')\) is therefore invertible. Then \(a / p = (a / e) \cdot (e / p)\) , whence

\[
a / p = \varepsilon (a) \varepsilon (p) ^ {*}\tag{1}
\]

for \(a \in \mathbf{E}\) and \(p \in \mathbf{S}'\) , the monoid \(\mathbf{E}_{\mathbf{S}}\) is therefore generated by \(\varepsilon(\mathbf{E}) \cup \varepsilon(\mathbf{S})^*\) .

PROPOSITION 6. The notation is that of Definition 7 and \(\varepsilon\) denotes the canonical homomorphism of E into \(\mathbf{E}_{\mathrm{S}}\) .

\begin{enumerate}
\def\labelenumi{(\roman{enumi})}
\item
  Let \(a\) and \(b\) be in \(\mathbf{E}\) ; in order that \(\varepsilon(a) = \varepsilon(b)\) , it is necessary and sufficient that there exist \(s \in \mathbf{S}'\) with \(sa = sb\) .
\item
  For \(\varepsilon\) to be injective it is necessary and sufficient that every element of S be cancellable.
\item
  For \(\varepsilon\) to be bijective it is necessary and sufficient that every element of \(S\) be invertible.
\end{enumerate}

Assertion (i) is clear and shows that \(\varepsilon\) is injective if and only if every element of \(S'\) is cancellable; but since the set of cancellable elements of \(E\) is a submonoid of \(E\) (no. 2, Proposition 2), it amounts to the same to say that every element of \(S\) is cancellable.

If \(\varepsilon\) is bijective every element of S is invertible, for \(\varepsilon(S)\) is composed of invertible elements of \(E_{S}\) . Conversely, suppose every element of S is invertible; then every element of \(S'\) is invertible (no. 3, Corollary 2 to Proposition 4) and hence cancellable. Then \(\varepsilon\) is injective by (ii) and \(a/p = \varepsilon(a.p^{*})\) by (1), hence \(\varepsilon\) is surjective.

THEOREM 1. Let E be a commutative monoid, S a subset of E, \(E_{S}\) the monoid of fractions associated with S and \(\varepsilon: E \to E_{S}\) the canonical homomorphism. Further let f be a homomorphism of E into a monoid F (not necessarily commutative) such that every element of \(f(S)\) is invertible in F. There exists one and only one homomorphism \(\bar{f}\) of \(E_{S}\) into F such that \(f = \bar{f} \circ \varepsilon\) .

If \(\bar{f}\) is a homomorphism of \(\mathbf{E}_{\mathbf{S}}\) into \(\mathbf{F}\) such that \(f = \bar{f} \circ \varepsilon\) , then

\[
\bar {f} (a | p) = \bar {f} (\varepsilon (a) \varepsilon (p) ^ {*}) = \bar {f} (\varepsilon (a)) \bar {f} (\varepsilon (p)) ^ {*} = f (a) f (p) ^ {*}
\]

for \(a \in \mathbf{E}\) and \(p \in \mathbf{S}'\) , whence the uniqueness of \(\bar{f}\) .

Let \(g\) be the mapping of \(\mathbf{E} \times \mathbf{S}'\) into \(\mathbf{F}\) defined by \(g(a, p) = f(a).f(p)^*\) . We show that \(g\) is a homomorphism of \(\mathbf{E} \times \mathbf{S}'\) into \(\mathbf{F}\) . First of all,

\[
g (e, e) = f (e) f (e) ^ {*} = e.
\]

Let \((a, p)\) and \((a', p')\) be two elements of \(\mathbf{E} \times \mathbf{S}'\) ; as \(a'\) and \(p\) commute in \(\mathbf{E}\) , \(f(a')\) and \(f(p)\) commute in \(\mathbf{F}\) , whence \(f(a')f(p)^* = f(p)^*f(a')\) by no. 3, Proposition 5. Moreover \(f(p p')^* = f(p' p)^* = (f(p')f(p))^* = f(p)^*f(p')^*\) by no. 3, Proposition 4, whence

\[
\begin{array}{l} g (a a ^ {\prime}, p p ^ {\prime}) = f (a a ^ {\prime}) f (p p ^ {\prime}) ^ {*} = f (a) f (a ^ {\prime}) f (p) ^ {*} f (p ^ {\prime}) ^ {*} = f (a) f (p) ^ {*} f (a ^ {\prime}) f (p ^ {\prime}) ^ {*} \\ = g (a, p) g (a ^ {\prime}, p ^ {\prime}). \end{array}
\]

We show that \(g\) is compatible with the equivalence relation \(\mathbf{R}\) on \(\mathbf{E} \times \mathbf{S}'\) : if \((a, p)\) and \((a', p')\) are congruent mod. \(\mathbf{R}\) , there exists \(s \in \mathbf{S}'\) with \(spa' = sap'\) , whence \(f(s)f(p)f(a') = f(s)f(a)f(p')\) . As \(f(s)\) is invertible, it follows that \(f(p)f(a') = f(a)f(p')\) and then by left multiplication by \(f(p)^*\) and right multiplication by \(f(p')^*\)

\[
g (a ^ {\prime}, p ^ {\prime}) = f (a ^ {\prime}) f (p ^ {\prime}) ^ {*} = f (p) ^ {*} f (a) = f (a) f (p) ^ {*} = g (a, p).
\]

Hence there exists a homomorphism \(\bar{f}\) of \(\mathbf{E}_{\mathbf{S}}\) into \(\mathbf{F}\) such that \(\bar{f}(a/p) = g(a,p)\) , whence \(\bar{f}(\varepsilon(a)) = \bar{f}(a/e) = f(a)f(e)^* = f(a)\) . Hence \(\bar{f} \circ \varepsilon = f\) .

COROLLARY. Let E and F be two commutative monoids, S and T subsets of E and F respectively, f a homomorphism of E into F such that \(f(\mathrm{S}) \subset \mathrm{T}\) and \(\varepsilon: E \to E_{S}\) , \(\eta: F \to F_{T}\) the canonical homomorphisms. There exists one and only one homomorphism \(g: E_{S} \to F_{T}\) such that \(g \circ \varepsilon = \eta \circ f\) .

The homomorphism \(\eta \circ f\) of \(E\) into \(F_{T}\) maps every element of \(S\) to an invertible element of \(F_{T}\) .

Remarks. (1) Theorem 1 can also be expressed by saying that \((\mathbf{E}_{\mathbf{S}},\varepsilon)\) is the solution of the universal mapping problem for \(\mathbf{E}\) , relative to monoids, monoid homomorphisms and homomorphisms of \(\mathbf{E}\) into monoids which map the elements of \(\mathbf{S}\) to invertible elements (Set Theory, IV, § 3, no. 1). It follows (loc. cit.) that every other solution of this problem is isomorphic in a unique way to \((\mathbf{E}_{\mathbf{S}},\varepsilon)\) .

\begin{enumerate}
\def\labelenumi{(\arabic{enumi})}
\setcounter{enumi}{1}
\tightlist
\item
  For the existence of a solution to the above universal mapping problem it is unnecessary to assume that the monoid E is commutative, as follows from Set Theory, IV, § 3, no. 2 (cf.~Exercise 17).
\end{enumerate}

We mention two important special cases of monoids of fractions.

\begin{enumerate}
\def\labelenumi{(\alph{enumi})}
\item
  Let \(\overline{\mathbf{E}} = \mathbf{E}_{\mathbb{E}}\) . As the monoid \(\overline{\mathbf{E}}\) is generated by the set \(\varepsilon(\mathbf{E}) \cup \varepsilon(\mathbf{E})^*\) , which is composed of invertible elements, every element of \(\overline{\mathbf{E}}\) is invertible (no. 3, Corollary 2 to Proposition 4). In other words, \(\overline{\mathbf{E}}\) is a commutative group. Moreover, by Theorem 1 every homomorphism \(f\) of \(\mathbf{E}\) into a group \(\mathbf{G}\) can be uniquely factorized in the form \(f = \bar{f} \circ \varepsilon\) , where \(\bar{f}: \overline{\mathbf{E}} \to \mathbf{G}\) is a homomorphism. \(\overline{\mathbf{E}}\) is called the group of fractions of \(\mathbf{E}\) (or group of differences of \(\mathbf{E}\) in the case of additive notation).
\item
  Let \(\Phi = \mathrm{E}_{\Sigma}\) , where \(\Sigma\) consists of the cancellable elements of \(\mathbf{E}\) . By Proposition 6, (ii), the canonical homomorphism of \(\mathbf{E}\) into \(\Phi\) is injective; it will be profitable to identify \(\mathbf{E}\) with its image in \(\Phi\) . Hence \(\mathbf{E}\) is a submonoid of \(\Phi\) , every cancellable element of \(\mathbf{E}\) has an inverse in \(\Phi\) and every element of \(\Phi\) is of the form \(a/p = a.p^{*}\) with \(a \in \mathbf{E}\) and \(p \in \Sigma\) ; then \(a/p = a'/p'\) if and only if \(ap' = pa'\) . It is easily seen that the invertible elements of \(\Phi\) are the fractions \(a/p\) , where \(a\) and \(p\) are cancellable, and \(p/a\) is the inverse of \(a/p\) .
\end{enumerate}

Now let S be a set of cancellable elements of E and S' be the submonoid of E generated by S. If \(a/p\) and \(a'/p'\) are two elements of \(\mathbf{E}_{\mathbf{S}}\) , then \(a/p = a'/p'\) if and only if \(ap' = pa'\) (for \(sap' = spa'\) implies \(ap' = pa'\) for all \(s \in \mathbf{S}'\) ). \(\mathbf{E}_{\mathbf{S}}\) may therefore be identified with the submonoid of \(\Phi\) generated by \(\mathbf{E} \cup \mathbf{S}^*\) .

If every element of E is cancellable, then \(\Phi = \bar{E}\) and E is a submonoid of the commutative group \(\Phi\) . Conversely, if E is isomorphic to a submonoid of a group, every element of E is cancellable.

